% archon:physics
% archon:covers ArchonPhysics/HittingTime.lean

\chapter{Generic positive-time hitting and persistence interfaces}
\label{ch:ArchonPhysics_HittingTime}

This clean-room interface accepts an arbitrary event predicate on extended
nonnegative time.  It is the order-theoretic layer below a concrete
equipartition observable: it never asserts that an event occurs, that an
infimum is attained, or that the predicate is measurable or continuous.  The
strict positivity guard excludes a zero-time observation.  These choices match
the operational definitions in the project roadmap, Section~6, while keeping
the library reusable for any event predicate.

\begin{definition}\leanok
[Strictly positive hitting times]
\label{def:hitting:times}
\lean{ArchonPhysics.HittingTime.hittingTimes}
For an event $A:\overline{\mathbb R}_{\geq0}\to\mathrm{Prop}$, let
\[
  H(A)=\{t\mid 0<t\mathbin{\wedge}A(t)\}.
\]
\end{definition}
\begin{proof}

\leanok
This is the set comprehension of positive times at which the supplied event
holds.  It is a definition and makes no existence assertion.
\end{proof}

\begin{definition}\leanok
[First hitting time]
\label{def:hitting:first-time}
\lean{ArchonPhysics.HittingTime.firstHittingTime}
\uses{def:hitting:times}
Define $\tau(A)=\inf H(A)$ in the complete order of extended nonnegative
reals.  Thus an empty event set has value $\top$.
\end{definition}
\begin{proof}

\leanok
Take the complete-lattice infimum of the preceding set.  The convention for
the infimum of an empty subset is the top element, so no separate convention
or attainment hypothesis is required.
\end{proof}

\begin{theorem}\leanok
[Hitting-time defining formulae]
\label{thm:hitting:spec}
\lean{ArchonPhysics.HittingTime.hitting_spec}
\uses{def:hitting:times, def:hitting:first-time}
The hitting-time set is exactly $H(A)$ and $\tau(A)=\inf H(A)$.
\end{theorem}
\begin{proof}

\leanok
Both conjuncts follow by unfolding the two transparent definitions.
\end{proof}

\begin{theorem}\leanok
[The empty event has infinite first hit]
\label{thm:hitting:empty}
\lean{ArchonPhysics.HittingTime.firstHittingTime_empty}
\uses{def:hitting:times, def:hitting:first-time}
For the identically false event, $\tau(A)=\top$.
\end{theorem}
\begin{proof}

\leanok
The positive hitting-time set is empty.  Its infimum in the extended
nonnegative reals is $\top$.
\end{proof}

\begin{theorem}\leanok
[Every hit bounds the first hit]
\label{thm:hitting:le-of-mem}
\lean{ArchonPhysics.HittingTime.firstHittingTime_le_of_mem}
\uses{def:hitting:times, def:hitting:first-time}
If $t\in H(A)$, then $\tau(A)\leq t$.
\end{theorem}
\begin{proof}

\leanok
An infimum is below every member of the set of which it is the infimum.
\end{proof}

\begin{theorem}\leanok
[Event inclusion reverses first-hit order]
\label{thm:hitting:mono}
\lean{ArchonPhysics.HittingTime.firstHittingTime_mono}
\uses{def:hitting:times, def:hitting:first-time}
If $A(t)$ implies $B(t)$ for every $t$, then $\tau(B)\leq\tau(A)$.
\end{theorem}
\begin{proof}

\leanok
The implication gives $H(A)\subseteq H(B)$.  Taking the infimum of a larger
set can only decrease it, which proves the stated reverse order.
\end{proof}

\begin{theorem}\leanok
[No positive event strictly before the infimum]
\label{thm:hitting:not-event-before}
\lean{ArchonPhysics.HittingTime.not_event_before_firstHittingTime}
\uses{def:hitting:times, def:hitting:first-time, thm:hitting:le-of-mem}
If $0<t<\tau(A)$, then $A(t)$ is false.
\end{theorem}
\begin{proof}

\leanok
If $A(t)$ held, strict positivity would put $t$ in $H(A)$; the preceding
bound would give $\tau(A)\leq t$, contradicting $t<\tau(A)$.  This is only a
pre-hit exclusion and does not claim that $A(\tau(A))$ holds.
\end{proof}

\begin{definition}\leanok
[Closed-interval persistence]
\label{def:hitting:persists-for}
\lean{ArchonPhysics.HittingTime.PersistsFor}
Say that $A$ persists from $t$ for duration $L$ when it holds for every
$s\in[t,t+L]$.
\end{definition}
\begin{proof}

\leanok
This is the universal predicate over the closed extended-time interval.  It
does not require any regularity of $A$.
\end{proof}

\begin{definition}\leanok
[Positive persistence starts]
\label{def:hitting:persistence-times}
\lean{ArchonPhysics.HittingTime.persistenceTimes}
\uses{def:hitting:persists-for}
Let $P(A,L)=\{u\mid0<u\mathbin{\wedge}\operatorname{PersistsFor}(A,u,L)\}$.
\end{definition}
\begin{proof}

\leanok
This is the same strict-positive-time guard applied to the duration-indexed
persistence predicate.
\end{proof}

\begin{definition}\leanok
[Persistence-time infimum]
\label{def:hitting:persistence-time}
\lean{ArchonPhysics.HittingTime.persistenceTime}
\uses{def:hitting:persistence-times}
Define $\tau_{\rm pers}(A,L)=\inf P(A,L)$, again allowing $\top$ when no
such start exists.
\end{definition}
\begin{proof}

\leanok
Take the extended-nonnegative infimum of the preceding set; completeness also
covers the empty case without an occurrence or attainment assumption.
\end{proof}

\begin{theorem}\leanok
[Persistence defining formulae]
\label{thm:hitting:persistence-spec}
\lean{ArchonPhysics.HittingTime.persistence_spec}
\uses{def:hitting:persists-for, def:hitting:persistence-times,
      def:hitting:persistence-time}
Persistence is precisely universal holding on $[t,t+L]$; its start set and
time are respectively $P(A,L)$ and its infimum.
\end{theorem}
\begin{proof}

\leanok
Unfold the three definitions.  The first conjunct exposes the closed-interval
predicate, and the remaining conjuncts expose the set comprehension and its
infimum.
\end{proof}
